\section{Background}

Self-Refinement improves an Initial Candidate through one or more subsequent model calls.
Self-Refine established a general pattern in which the same large language model (LLM) generates an output, produces natural-language feedback, and revises the output without additional training \cite{madaan2023selfrefine}.
In code generation, Self-Debugging prompts the model to explain and revise its own code and, when unit tests are available, also uses execution results as feedback \cite{chen2024selfdebug}.
We use Self-Refinement in a narrower sense: the same model reviews and revises its Self-Generated Initial Candidate using only the task specification and self-generated intermediate artifacts, while tests are reserved for post-hoc evaluation and no Execution Feedback is provided to the model.

\subsection{Challenges and Risks of Code Self-Refinement}

An additional refinement call does not necessarily produce a better program.
Olausson et al. compare self-repair with independent resampling at an equivalent program-sampling budget and find modest or absent gains in several settings; stronger-model and human feedback improve repair, indicating that feedback quality is a key bottleneck \cite{olausson2024selfrepair}.
For locally run 1-3B models, Feedback Over Form finds that execution feedback and early stopping drive improvements, while its constrained topology search provides no clear gain over a simple refinement loop \cite{mcandrews2026feedback}.
Its exploratory text-only pipeline experiments show no improvement, but the study treats them as secondary, less systematic evidence \cite{mcandrews2026feedback}.
Although both studies examine repair with execution feedback, they motivate evaluating whether the benefits of additional refinement calls justify their cost.

Refinement can also damage an initially correct program.
CoCoS explicitly reports incorrect-to-correct and correct-to-incorrect transitions and trains small models to maintain correct outputs while improving incorrect ones \cite{cho-etal-2025-self}.
This transition-based view is close to our analysis, but CoCoS trains a correction-preserving model rather than comparing inference-time protocols over the same Initial Candidate.
CodeChat-Eval provides complementary evidence by finding statistically significant declines in Functional Correctness across ten-turn sessions of functionality-preserving follow-up instructions \cite{guo2026codechateval}.
Its revisions follow curated developer-preferred instructions rather than self-generated critique, but the observed regressions demonstrate that more editing is not inherently beneficial.

A refinement system must therefore decide not only how to revise a candidate, but also whether the candidate should be revised.
Jin and Chen show that, when judging requirement conformance without executing tests, LLMs frequently reject correct implementations and justify these decisions with invented requirements, excessive edge-case concerns, or unsupported claims \cite{jin2026reliable}.
Such overcorrection may trigger unnecessary refinement, although a false review verdict does not by itself establish a Correct-to-Incorrect Regression in the revised code.
Revisit Self-Debugging finds that post-execution repair can be misled by bias in self-generated tests, showing that apparently informative test feedback may still provide an unreliable basis for repair \cite{chen-etal-2025-revisit}.
These findings motivate paired analysis of Repair, Correct-to-Incorrect Regression, and exact preservation instead of relying only on final pass rate.

\subsection{Methods for Improving Code Self-Refinement}

One family of methods improves refinement through explicit critique or reflection.
RCO trains a critic with a critique-utility reward derived from the improvement of an actor's revised response, avoiding direct critique-preference assessment \cite{yu-etal-2025-training}.
SCOPE adapts a prover model into a structured critic that produces subgoals, gap analysis, and a robustness checklist to guide code correction \cite{zhang2026scope}.
Broader reflection systems include Reflexion, which stores verbal reflections on task feedback in episodic memory \cite{shinn2023reflexion}, and INDICT, which coordinates safety- and helpfulness-oriented critics equipped with external knowledge and tools \cite{le2024indict}.
These methods support the idea that an explicit critique or reflection artifact can guide subsequent decisions, but they use trained critics, external signals and tools, or integrated agent loops rather than isolating the effect of adding Critique and Planning to a fixed model.

Other methods learn correction behavior or structure pre-execution review.
SCoRe trains self-correction with two-stage, multi-turn online reinforcement learning under the model's own distribution of self-generated correction traces \cite{kumar2025score}, whereas CoCoS adapts multi-turn reinforcement learning to small models through accumulated and progressive rewards \cite{cho-etal-2025-self}.
ReflexiCoder internalizes initial generation, bug- and optimization-aware reflection, and self-correction in the model weights, avoiding an execution engine at inference time \cite{jiang-etal-2026-reflexicoder}.
ReVISE uses a preference-learning curriculum to train self-verification and reasoning correction sequentially \cite{pmlr-v267-lee25ab}, while RubricRefine generates task- and tool-registry-specific rubrics to check semantic contracts and repair tool-use code before execution \cite{levine2026rubricrefine}.
These approaches show that correction and verification can be learned or structured, but their integrated procedures do not expose the incremental correctness and token effects of separate inference-time Refinement Stages.
Code Reffix uses oracle reflections and a dual-task protocol to decouple the evaluation of reflection generation from reflection-guided repair \cite{di-etal-2026-code}, and CodeReviewQA decomposes review comprehension into change-type recognition, change localization, and solution identification \cite{lin-etal-2025-codereviewqa}.
The COAST study evaluates bug localization, identification, repair, and recognition through DEBUGEVAL and uses communicative agents to synthesize debugging data for training \cite{yang-etal-2025-coast}.
These decompositions provide indirect support for Stage Composition, while focusing on capability evaluation or training rather than Shared-Initial protocol comparison.

A second family improves refinement with Execution Feedback.
Self-Edit supplies example-test execution results to a fault-aware editor \cite{zhang-etal-2023-self}, while CYCLE trains models to refine faulty generations from available feedback, including test-suite execution results \cite{ding2024cycle}.
LeDex trains on explanation-and-refinement trajectories filtered through execution verification \cite{jiang2024ledex}, and OpenCodeInterpreter supports iterative code generation and refinement using execution and human feedback \cite{zheng-etal-2024-opencodeinterpreter}.
More detailed signals are used by LDB, which tracks and verifies runtime states at the basic-block level \cite{zhong-etal-2024-debug}, and by TraceCoder, which uses runtime traces for diagnosis, repair, and rollback \cite{huang2026tracecoder}.
CodeT instead selects among generated solutions using generated tests and dual execution agreement \cite{chen2023codet}, whereas LEVER reranks generated programs with an execution-aware verifier \cite{pmlr-v202-ni23b}; neither method revises a single Initial Candidate.
These methods demonstrate the value of behavioral evidence for debugging, but their results cannot determine how well a model refines code without execution feedback or external verification.

Multi-agent and repository-level systems place refinement within broader workflows: AgentCoder separates programming, test design, and test execution \cite{huang2023agentcoder}; MapCoder separates example recall, planning, code generation, and debugging \cite{islam-etal-2024-mapcoder}; RepairAgent autonomously selects tools for program repair \cite{bouzenia2025repairagent}; and Agentless reduces repository repair to localization, repair, and patch validation \cite{xia2024agentless}.
Their stage and resource tradeoffs therefore include execution or repository interaction that is absent from our fixed-candidate setting.
Adjacent work also includes InterCode, which formalizes interactive coding as reinforcement-learning environments where execution feedback is observed after code actions \cite{yang2023intercode}, and RefineCoder, which iteratively retrains code models on data refined from self-generated code and external critiques \cite{zhou2025refinecoder}.
SemCoder trains models to reason about statement-level execution effects and program behavior \cite{ding2024semcoder}, while CRITIC validates and revises outputs using feedback from external tools \cite{gou2024critic}.
EffiLearner feeds execution-time and memory profiles back to the model to optimize generated code \cite{huang2024effilearner}, whereas FunCoder recursively decomposes tasks into subfunctions and uses behavior-based consensus to select implementations \cite{chen2024funcoder}.
ReChisel refines Chisel code using compiler and simulation feedback \cite{niu2025rechisel}, while RethinkMCTS uses code-execution feedback to repair erroneous thoughts during Monte Carlo tree search \cite{li-etal-2025-rethinkmcts}.
SWE-agent provides a custom agent-computer interface for repository navigation, code editing, and program execution \cite{yang2024sweagent}, while MdEval benchmarks program repair, bug localization, and bug identification across 20 programming languages \cite{liu-etal-2026-mdeval}.
Other work trains a lightweight detector on internal model representations to identify line-level code hallucinations \cite{andriushchenko-etal-2026-efficient}, and CodeRise combines progressive curriculum generation with multi-turn self-debugging for ultra-low-resource programming languages \cite{wen-etal-2026-coderise}.
These adjacent directions broaden the literature but do not directly compare inference-time revisions of the same candidate.

Existing studies therefore establish that refinement may repair or regress code and that critique, training, selective verification, and Execution Feedback can improve parts of the process, but evidence remains limited for inference-only Self-Refinement by locally deployable open-weight instruct models under a strict no-feedback condition.
In particular, none of the reviewed studies jointly uses a Shared-Initial Design to compare Revision alone, Critique followed by Revision, and Critique followed by Planning and Revision as cumulative Stage Compositions; measures Repair and Correct-to-Incorrect Regression for the same candidates; evaluates whether a standalone Decision prevents regressions or misses repairs; and accounts for the token consumption of each stage.
This study addresses that gap by applying the same model to the same Initial Candidate across all protocols, incrementally adding Critique and Planning before Revision, and withholding execution feedback.
The resulting comparison characterizes the correctness and token-cost tradeoffs of protocol structure and Decision-Conditioned Refinement before the following section formalizes the experimental design.
